\section{Pretrained-model aggregation details}
\label{app:lment-frequency}

The source analysis uses seed 42 and an object-coverage-balanced ordering to
choose eight LRE estimation examples per relation. These examples are excluded
from evaluation and supply the in-context demonstrations for evaluation prompts.
The filter retains a sample when the model's decoded top-1 token, stripped of
whitespace and lowercased, is a nonempty prefix of the gold label. This is a
first-token criterion and does not certify the full answer. Candidate subject
states are residual states after blocks $\{0,3,6,9,12,15\}$, indexed from zero.
The archived scale sweep uses $\{1,2,3,4,5,6,8,10\}$ and maximizes the macro mean
of each relation's best-layer faithfulness. It selects $\beta=4$. Layer and scale
selection both use the reported evaluation outcomes; the estimation split does
not supply an independent model-selection test set.

For each relation and frequency axis, let $q_{.25}$ and $q_{.75}$ be empirical
quartiles computed by linear interpolation between adjacent sorted observations.
The lower group contains every fact with $f\leq q_{.25}$ and the upper group every
fact with $f\geq q_{.75}$. We require $q_{.25}<q_{.75}$ so that these groups are
disjoint. Discrete counts can produce groups larger than one quarter of the
relation; we retain all tied observations. The difference is
\begin{equation}
  \Delta_r=100\left(
    \frac{\sum_{i:f_i\geq q_{.75}} y_i}{|\{i:f_i\geq q_{.75}\}|}
    -\frac{\sum_{i:f_i\leq q_{.25}} y_i}{|\{i:f_i\leq q_{.25}\}|}
  \right),
\end{equation}
where $y_i$ is the binary faithfulness outcome. We average $\Delta_r$ with equal
weight per eligible relation. The pooled Spearman correlations instead give
every evaluated fact equal weight and use average ranks for ties. No independence
test or uncertainty estimate over model seeds is implied by these summaries.

Ten relations, totaling 505 evaluated facts, have only one distinct model-predicted
answer token. All have perfect LRE faithfulness. In particular, all 203
person-employer answers and all 154 person-educated-at answers begin with
\texttt{University}. A constant decoder would match the model on these sets.
We retain them in the all-fact reproduction of the source curves, but exclude
them from the within-relation comparison. Removing all ten leaves 4,883 facts
across 48 relations with $46.0\%$ pooled faithfulness. The pooled rank correlations
then become $0.073$, $0.008$, and $-0.033$ for subject, subject--object, and
subject--relation frequency. The change in the subject--object correlation shows
why the pooled statistic alone is insufficient. Token imbalance within the
remaining relations can still affect the comparisons.

The presence of decreases is not specific to the ten-fact tail threshold. With
at least five facts per tail, the numbers of decreases are $15/34$, $14/33$, and
$11/26$ eligible relations for the three axes. With at least twenty, they are
$9/20$, $9/20$, and $7/14$. These thresholds were checked during the aggregation;
they were not preregistered. Eligibility depends on frequency variation and sample
size, so these proportions describe the retained relations rather than an
estimate for all natural-language relations.

The paper's \texttt{analyze\_lment\_frequency.py} script rebuilds the figures and
tables from \texttt{data/lment\_frequency/facts.csv}. The snapshot joins predictions
to document counts by relation, subject, and object, with a one-to-one merge and
no missing counts. Its provenance manifest records input hashes and the original
evaluation and counting configurations. We checked every exact-frequency count
against all 174 raw curves in the original browser payload and verified that
estimation facts do not overlap evaluation facts. The analysis files include
per-bin denominators, every relation's quartile cutoffs and scores, and the tail
size sensitivity results. No model was retrained or re-evaluated for this
aggregation. The pretrained evidence is limited to this checkpoint; the synthetic
interventions remain the causal tests of the frequency-only explanation.
